\section{引言}
\label{sec:intro}

时变再生数估计已成为传染病监测的常规产出。此类估计曾用于追踪西非埃博拉病毒病的传播\citep{nouvellet2018simple}、评估针对 COVID-19 的非药物干预\citep{flaxman2020estimating}，以及描述 SARS-CoV-2 的早期传播\citep{kucharski2020early}。目前许多国家已常规发布再生数估计\citep{huisman2022estimation}，其估计方法也有实践指南可循\citep{gostic2020practical, steyn2025primer}。再生数估计对态势感知的价值，取决于依据监测系统实际收集的数据能够把它估计到何种精度。

瞬时再生数 $R_t$ 定义为：若 $t$ 时刻的条件保持不变，在该时刻被感染的个体所产生的二代感染的期望数\citep{fraser2007estimating}。$R_t$ 通常借助更新方程从发病数估计，更新方程按代际间隔（原发病例被感染到其二代病例被感染之间的时间）对既往发病数加权\citep{wallinga2004different, cori2013new, thompson2019improved}。在代际间隔给定时，此类估计的后验不确定性随估计窗口内病例数的增加而缩小\citep{cori2013new}。Parag 等人以闭式刻画了发病数所含关于 $R_t$ 的信息量，其中包括随时间变化的报告比例造成的信息损失，并推导了这一信息量对及时决策的限制\citep{parag2021improved, parag2022quantifying, parag2025asymmetric}。已有扩展处理了时间聚合的计数\citep{nash2023estimating, ogigittins2024simulation}，以及感染的再生数与延迟观测序列的表观再生数之间的差异\citep{eales2024differences}。

上述精度结果采用同一基准，即误差相对于所估时期 $R_t$ 的实现值来度量，并随病例数增加而趋于零。$R_t$ 估计的许多用途则指向跨期持续存在的量，例如一个流行季的典型传播水平或下一周的期望值。政策评估同样要对天气、学校日历和接触模式各不相同的若干周取平均\citep{dalziel2018urbanization}。$R_t$ 在时期之间的变异为同一时间人群中的全部传播链所共有。种群生态学把这类环境随机性与源于个体独立结局的种群统计随机性区分开来\citep{engen1998demographic, rohani2002determinism, lande2003stochastic}，机制性流行病模型也已同时纳入这两种随机性\citep{breto2009time}。每个个体分摊的种群统计方差随种群规模增大而下降，环境方差则保持不变；两种贡献相等时的种群规模是灭绝风险的标准概括量\citep{lande2003stochastic}。过度离散的后代分布\citep{lloydsmith2005superspreading, endo2020estimating}与不完全报告都归入这一划分中的种群统计一侧。

监测数据还引入第三个变异来源。进入计数的感染比例随检测实践、就医行为和医疗机构报告完整度而随时间变化\citep{azmon2014estimating, stoner2019hierarchical, parag2022fundamental}。在美国，医院向美国国家医疗安全网络（NHSN）报告流感住院人数在 2024 年 4 月之前为强制要求，2024 年 5 月至 10 月改为自愿报告，自 2024 年 11 月起恢复强制\citep{cdc2026nhsn}。此类住院序列越来越多地用于预测和实时评估\citep{mathis2024evaluation}。在一周之内持续存在的报告完整度变异作用于该周的每一个病例，因此其作用方式与环境变异相同，有别于抽样噪声。

监测的样本量计算通常假定病例越多，精度按比例越高\citep{cochran1977sampling}；关于再生数估计的分析则集中于发病数据对逐期取值的信息量\citep{parag2022quantifying, bhatt2023semi}。据我们所知，生态学的方差分解尚未被引入再生数监测的设计；引入时需要纳入过度离散的后代分布、完整度波动的二项报告，以及分摊到多个时期的预算。实证检验还面临一个测量问题。由住院数计算的周增长统计量属于表观量，其方差取决于由感染到入院的延迟以及各统计量在一周之内汇总增长的方式，因此其截距不能直接解读为传播方差。同一日历周的观测共同受到全国性因素的影响，把它们视为相互独立的区间估计会过窄\citep{cameron2011robust}。

本文把生态学的方差分解推广到 $R_t$ 存在环境变异、后代服从负二项分布、报告服从二项分布时基于病例数的再生数估计，并推导环境贡献与抽样贡献相等时的平衡规模。在此基础上，分析这一规模对报告异质性的含义，以及固定预算如何在每期病例数与时期数之间分配。随后，利用 NHSN 的每周流感住院数据，以两种表观周统计量估计这一规模，并通过数据匹配模拟的两步反演，把它们的截距映射为潜伏环境方差。最后，把截距分解为全国周分量与州特有分量，并比较强制报告与自愿报告两种制度。
